% =========================
% L19.tex
% =========================

\chapter{L19 -- Realizzazione in forma di stato di sistemi MIMO, zeri di trasmissione e stabilità}
\label{chap:L19}

\section{Dalla matrice di trasferimento alla forma di stato}

In questa lezione affrontiamo tre temi collegati tra loro:
\begin{enumerate}
    \item come leggere la struttura di un sistema \textbf{MIMO} a partire dalla sua matrice di trasferimento;
    \item come individuare i \textbf{poli}, gli \textbf{zeri di trasmissione} e l'\textbf{ordine minimo} della realizzazione;
    \item come distinguere \textbf{stabilità interna} e \textbf{stabilità esterna (BIBO)}.
\end{enumerate}

Partiamo da una matrice di trasferimento razionale
\[
G(s)\in \mathbb{R}(s)^{\ell\times m},
\]
dove \(m\) è il numero di ingressi e \(\ell\) il numero di uscite.

L'obiettivo finale è sempre quello di costruire una realizzazione in forma di stato
\[
\begin{cases}
\dot x = Ax + Bu,\\
y = Cx + Du,
\end{cases}
\qquad
B\in\mathbb{R}^{n\times m},\quad
C\in\mathbb{R}^{\ell\times n},
\]
che rappresenti il sistema dato.

Nel caso MIMO, però, non basta guardare una singola funzione di trasferimento: la struttura del sistema è contenuta nell'intera matrice \(G(s)\). Per questo diventa naturale introdurre la \textbf{forma canonica di Smith} per matrici polinomiali e, di conseguenza, la \textbf{forma di Smith--McMillan} per matrici razionali.

\bigskip

\section{Forma canonica di Smith per matrici polinomiali}

Sia
\[
P(s)\in \mathbb{R}[s]^{\ell\times m}
\]
una matrice a coefficienti polinomiali. Allora esistono due matrici unimodulari \(U_1(s)\) e \(U_2(s)\) tali che
\[
P(s)=U_1(s)\,\Lambda(s)\,U_2(s),
\]
dove \(\Lambda(s)\) è diagonale della forma
\[
\Lambda(s)=
\operatorname{diag}\bigl(\delta_1(s),\delta_2(s),\dots,\delta_r(s),0,\dots,0\bigr),
\]
con:
\begin{itemize}
    \item \(r=\operatorname{rk}P(s)\) (rango normale della matrice polinomiale);
    \item ogni \(\delta_i(s)\) monico;
    \item \(\delta_i(s)\mid \delta_{i+1}(s)\) per ogni \(i\).
\end{itemize}

Questa è la \textbf{forma canonica di Smith}. I polinomi \(\delta_i(s)\) sono i \textbf{fattori invarianti} della matrice \(P(s)\).

\medskip

\noindent
\textbf{Osservazione importante.}
La forma di Smith si applica a una \emph{matrice polinomiale}, non direttamente a una matrice razionale. Per questo, quando partiamo da \(G(s)\), prima la riscriviamo come rapporto tra una matrice polinomiale e un denominatore comune.

\bigskip

\section{Dalla forma di Smith alla forma di Smith--McMillan}

Supponiamo che la matrice di trasferimento possa essere scritta come
\[
G(s)=\frac{1}{d(s)}\,P(s),
\]
dove:
\begin{itemize}
    \item \(d(s)\) è un polinomio scalare comune ai denominatori;
    \item \(P(s)\in\mathbb{R}[s]^{\ell\times m}\).
\end{itemize}

Se per \(P(s)\) calcoliamo la forma di Smith
\[
P(s)=U_1(s)\,\Lambda(s)\,U_2(s),
\]
allora
\[
G(s)=U_1(s)\,\Pi(s)\,U_2(s),
\qquad
\Pi(s)=\frac{\Lambda(s)}{d(s)}.
\]

Esplicitamente,
\[
\Pi(s)=
\operatorname{diag}\!\left(
\frac{\delta_1(s)}{d(s)},
\frac{\delta_2(s)}{d(s)},
\dots,
\frac{\delta_r(s)}{d(s)},
0,\dots,0
\right).
\]

Ora ogni rapporto \(\delta_i(s)/d(s)\) va semplificato. Dopo aver eseguito \emph{tutte} le cancellazioni possibili, si ottiene la \textbf{forma di Smith--McMillan}
\[
\Pi_{SM}(s)=
\operatorname{diag}\!\left(
\frac{\alpha_1(s)}{\beta_1(s)},
\dots,
\frac{\alpha_r(s)}{\beta_r(s)},
0,\dots,0
\right),
\]
dove:
\begin{itemize}
    \item i polinomi \(\alpha_i(s)\) e \(\beta_i(s)\) sono coprimi;
    \item i \(\beta_i(s)\) raccolgono i poli strutturali del sistema;
    \item gli \(\alpha_i(s)\) raccolgono gli zeri strutturali, cioè gli \textbf{zeri di trasmissione}.
\end{itemize}

\medskip

\noindent
\textbf{Informazioni che si leggono dalla forma di Smith--McMillan.}
\begin{itemize}
    \item I \textbf{poli} della matrice di trasferimento sono le radici dei \(\beta_i(s)\).
    \item Gli \textbf{zeri di trasmissione} sono le radici degli \(\alpha_i(s)\).
    \item L'\textbf{ordine minimo di una realizzazione} è la somma dei gradi dei denominatori dopo cancellazione:
    \[
    n_{\min}=\sum_{i=1}^r \deg \beta_i(s).
    \]
    Questo numero è il \textbf{grado di McMillan}.
\end{itemize}

\bigskip

\section{Esempio sulla forma di Smith--McMillan}

Consideriamo la matrice polinomiale
\[
P(s)=
\begin{bmatrix}
1 & -1\\
s^2+s-4 & 2s^2-s-8\\
s^2-4 & 2s^2-8
\end{bmatrix}.
\]

Dai calcoli della lezione precedente si sa che la sua forma di Smith è
\[
\Lambda(s)=
\begin{bmatrix}
1 & 0\\
0 & (s+2)(s-2)\\
0 & 0
\end{bmatrix}.
\]

Supponiamo ora di avere la matrice di trasferimento
\[
G(s)=\frac{1}{s^2+3s+2}\,P(s).
\]
Poiché
\[
s^2+3s+2=(s+1)(s+2),
\]
otteniamo
\[
\Pi(s)=\frac{\Lambda(s)}{s^2+3s+2}
=
\begin{bmatrix}
\dfrac{1}{(s+1)(s+2)} & 0\\[1.1em]
0 & \dfrac{(s+2)(s-2)}{(s+1)(s+2)}\\[1.1em]
0 & 0
\end{bmatrix}.
\]

Semplificando il secondo elemento diagonale:
\[
\Pi_{SM}(s)=
\begin{bmatrix}
\dfrac{1}{(s+1)(s+2)} & 0\\[1.1em]
0 & \dfrac{s-2}{s+1}\\[1.1em]
0 & 0
\end{bmatrix}.
\]

A questo punto possiamo leggere direttamente:

\subsection*{Poli}
I denominatori non banali sono
\[
(s+1)(s+2),\qquad s+1.
\]
Quindi i poli sono:
\[
s=-1 \quad \text{(doppio nel grado complessivo)},\qquad s=-2.
\]

Infatti il prodotto dei denominatori diagonali è
\[
(s+1)(s+2)\cdot (s+1)=(s+1)^2(s+2).
\]

\subsection*{Zeri di trasmissione}
I numeratori non banali sono
\[
1,\qquad s-2.
\]
Dunque compare un solo zero di trasmissione:
\[
s=2.
\]

\subsection*{Ordine minimo}
L'ordine minimo della realizzazione è
\[
n_{\min}=\deg\big((s+1)(s+2)\big)+\deg(s+1)=2+1=3.
\]

Questa è un'informazione molto importante:
\begin{quote}
\emph{anche se una qualunque realizzazione non minima può avere più stati, la forma minima del sistema ha dimensione \(3\).}
\end{quote}

\bigskip

\section{Interpretazione degli zeri di trasmissione}

Negli appunti compare la domanda fondamentale:

\begin{quote}
Dato un ingresso non nullo, può esistere una condizione iniziale tale che l'uscita sia identicamente nulla?
\end{quote}

Questa è proprio l'idea intuitiva di \textbf{zero di trasmissione}.

\subsection{Caso SISO continuo}

Consideriamo un sistema SISO in tempo continuo:
\[
\begin{cases}
\dot x = Ax + Bu,\\
y = Cx + Du.
\end{cases}
\]

Cerchiamo un ingresso del tipo
\[
u(t)=e^{\delta t}
\]
e una traiettoria dello stato della forma
\[
x(t)=x_0 e^{\delta t},
\]
con \(\delta\in\mathbb{C}\).

Sostituendo nell'equazione di stato:
\[
\delta x_0 e^{\delta t}=Ax_0 e^{\delta t}+B e^{\delta t}.
\]
Dividendo per \(e^{\delta t}\neq 0\),
\[
(\delta I-A)x_0=B.
\]

Se \(\delta\) non è autovalore di \(A\), possiamo invertire \(\delta I-A\) e scrivere
\[
x_0=(\delta I-A)^{-1}B.
\]

Allora l'uscita vale
\[
y(t)=C x(t)+D u(t)
= C(\delta I-A)^{-1}B\,e^{\delta t}+D e^{\delta t},
\]
cioè
\[
y(t)=\bigl(C(\delta I-A)^{-1}B+D\bigr)e^{\delta t}.
\]

Ma il termine tra parentesi è proprio
\[
G(\delta)=C(\delta I-A)^{-1}B+D.
\]

Quindi
\[
y(t)=G(\delta)e^{\delta t}.
\]

\medskip

\noindent
Se \(\delta\) è uno zero di \(G(s)\), allora
\[
G(\delta)=0
\quad\Longrightarrow\quad
y(t)\equiv 0,
\]
pur con ingresso non nullo.

\medskip

\noindent
\textbf{Interpretazione.}
Uno zero di trasmissione è un valore dell'esponenziale d'ingresso per cui è possibile scegliere opportunamente lo stato iniziale in modo da annullare completamente l'uscita.

\subsection{Generalizzazione MIMO}

Nel caso MIMO l'idea è analoga, ma l'ingresso ha una \emph{direzione}:
\[
u(t)=v e^{\delta t},\qquad v\neq 0.
\]
Si cercano quindi \(x_0\neq 0\) e \(v\neq 0\) tali che
\[
\begin{cases}
(\delta I-A)x_0-Bv=0,\\
Cx_0+Dv=0.
\end{cases}
\]

In forma compatta:
\[
\begin{bmatrix}
\delta I-A & -B\\
C & D
\end{bmatrix}
\begin{bmatrix}
x_0\\
v
\end{bmatrix}
=0.
\]

La matrice
\[
\mathcal{R}(\delta)=
\begin{bmatrix}
\delta I-A & -B\\
C & D
\end{bmatrix}
\]
è la \textbf{matrice di Rosenbrock}. I suoi valori \(\delta\) per cui il rango cala sono gli zeri di trasmissione del sistema.

\bigskip

\section{Stabilità interna}

La \textbf{stabilità interna} riguarda l'evoluzione dello \emph{stato}, quindi dipende dagli autovalori di \(A\), non direttamente dalla funzione di trasferimento.

\medskip

\noindent
Per il sistema autonomo
\[
\dot x = Ax,
\]
le componenti della soluzione sono costruite a partire dai modi associati agli autovalori di \(A\).

\subsection{Modi elementari nel caso continuo}

Se \(\lambda\in\mathbb{R}\) è autovalore reale di \(A\), i modi associati sono del tipo
\[
e^{\lambda t},\quad
t e^{\lambda t},\quad
t^2 e^{\lambda t},\quad \dots
\]
a seconda della dimensione dei blocchi di Jordan.

Se invece \(\lambda=\sigma\pm j\omega\) è complesso, compaiono modi del tipo
\[
e^{\sigma t}\cos(\omega t),\qquad
e^{\sigma t}\sin(\omega t),
\]
eventualmente moltiplicati per potenze di \(t\) nel caso di blocchi di Jordan non banali:
\[
t^k e^{\sigma t}\cos(\omega t),\qquad
t^k e^{\sigma t}\sin(\omega t).
\]

\subsection{Modi elementari nel caso discreto}

Per il sistema
\[
x(k+1)=Ax(k),
\]
i modi associati a un autovalore \(\lambda\) sono del tipo
\[
\lambda^k,\qquad
k\lambda^k,\qquad
k^2\lambda^k,\qquad \dots
\]
sempre a seconda della struttura di Jordan.

\subsection{Condizioni di stabilità interna asintotica}

\paragraph{Tempo continuo.}
Il sistema è asintoticamente stabile se e solo se
\[
\Re(\lambda_i)<0
\qquad \forall\, \lambda_i\in\sigma(A).
\]

\paragraph{Tempo discreto.}
Il sistema è asintoticamente stabile se e solo se
\[
|\lambda_i|<1
\qquad \forall\, \lambda_i\in\sigma(A).
\]

\medskip

\noindent
\textbf{Osservazione.}
Se un autovalore ha parte reale nulla (TC) o modulo unitario (TD), si possono avere oscillazioni persistenti: il sistema non diverge, ma non converge neppure a zero. Questo è esattamente ciò che succede nell'oscillatore armonico non smorzato.

\bigskip

\section{Stabilità esterna o BIBO}

La \textbf{stabilità esterna} riguarda invece il legame ingresso--uscita.

\medskip

\noindent
\textbf{Definizione (BIBO).}
Un sistema è \emph{Bounded Input Bounded Output} se, per ogni ingresso limitato, l'uscita corrispondente a stato iniziale nullo rimane limitata.

In simboli:
\[
\|u(t)\|\le M_u\ \forall t\ge 0
\quad\Longrightarrow\quad
\exists M_y>0:\ \|y(t)\|\le M_y\ \forall t\ge 0.
\]

Nel caso LTI, la stabilità BIBO si legge dai poli della funzione di trasferimento, non dagli autovalori nascosti di una realizzazione non minima.

\subsection{Condizione sui poli}

\paragraph{Tempo continuo.}
Un sistema razionale proprio è BIBO stabile se tutti i poli della funzione di trasferimento hanno parte reale strettamente negativa:
\[
\Re(p_i)<0.
\]

\paragraph{Tempo discreto.}
Un sistema razionale proprio è BIBO stabile se tutti i poli stanno strettamente dentro il cerchio unitario:
\[
|p_i|<1.
\]

\medskip

\noindent
\textbf{Punto chiave.}
I poli della funzione di trasferimento corrispondono solo agli autovalori di \(A\) che sono \textbf{sia raggiungibili sia osservabili}. Per questo:
\begin{itemize}
    \item la stabilità interna guarda \emph{tutti} gli autovalori di \(A\);
    \item la stabilità esterna guarda solo i modi che compaiono effettivamente in ingresso--uscita.
\end{itemize}

\bigskip

\section{Esempio: un sistema non BIBO stabile}

Consideriamo
\[
G(s)=\frac{1}{s},
\qquad
U(s)=\frac{1}{s}.
\]
Allora
\[
Y(s)=G(s)U(s)=\frac{1}{s^2}.
\]

Nel tempo, questo corrisponde a una rampa:
\[
y(t)=t.
\]
L'ingresso \(u(t)=1\) è limitato, ma l'uscita cresce senza limite. Dunque il sistema \(\frac{1}{s}\) \textbf{non è BIBO stabile}.

\medskip

Questo esempio chiarisce bene che, per avere BIBO stabilità, non basta che i modi non esplodano immediatamente: serve che l'uscita resti limitata per ogni ingresso limitato.

\bigskip

\section{Oscillatore armonico non smorzato}

Consideriamo il sistema
\[
\ddot y + y = u.
\]

Introduciamo le variabili di stato
\[
x_1=y,\qquad x_2=\dot y.
\]
Allora
\[
\begin{cases}
\dot x_1 = x_2,\\
\dot x_2 = -x_1 + u,
\end{cases}
\]
ossia
\[
\dot x =
\begin{bmatrix}
0 & 1\\
-1 & 0
\end{bmatrix}x+
\begin{bmatrix}
0\\
1
\end{bmatrix}u,
\qquad
y=
\begin{bmatrix}
1 & 0
\end{bmatrix}x.
\]

\subsection{Stabilità interna}

La matrice di stato è
\[
A=
\begin{bmatrix}
0 & 1\\
-1 & 0
\end{bmatrix},
\]
i cui autovalori sono
\[
\lambda_{1,2}=\pm j.
\]

Quindi il sistema autonomo non è asintoticamente stabile: i modi sono sinusoidali puri. Con \(u=0\), le traiettorie sono chiuse e non convergono all'origine.

Infatti:
\[
\frac{d}{dt}(x_1^2+x_2^2)
=2x_1\dot x_1+2x_2\dot x_2
=2x_1x_2+2x_2(-x_1)=0,
\]
quindi
\[
x_1^2(t)+x_2^2(t)=\text{costante}.
\]

Dunque le traiettorie stanno su circonferenze concentriche.

\begin{figure}[h!]
\centering
\begin{tikzpicture}[scale=1.1]
    % assi
    \draw[->] (-2.6,0) -- (2.6,0) node[right] {$x_1$};
    \draw[->] (0,-2.6) -- (0,2.6) node[above] {$x_2$};

    % orbite
    \draw[dashed] (0,0) circle (1.1);
    \draw (0,0) circle (2.0);

    % frecce orbita esterna (senso orario)
    \draw[->] (2.0,0) arc[start angle=0,end angle=-30,radius=2.0];
    \draw[->] (0,-2.0) arc[start angle=-90,end angle=-120,radius=2.0];
    \draw[->] (-2.0,0) arc[start angle=180,end angle=150,radius=2.0];
    \draw[->] (0,2.0) arc[start angle=90,end angle=60,radius=2.0];

    % frecce orbita interna
    \draw[->] (1.1,0) arc[start angle=0,end angle=-35,radius=1.1];
    \draw[->] (0,-1.1) arc[start angle=-90,end angle=-125,radius=1.1];
    \draw[->] (-1.1,0) arc[start angle=180,end angle=145,radius=1.1];
    \draw[->] (0,1.1) arc[start angle=90,end angle=55,radius=1.1];

    \node at (1.55,1.6) {\small \(x_1^2+x_2^2=c\)};
\end{tikzpicture}
\caption{Ritratto di fase dell'oscillatore armonico non smorzato.}
\end{figure}

\subsection{Stabilità esterna}

Calcoliamo la funzione di trasferimento:
\[
G(s)=C(sI-A)^{-1}B
=
\begin{bmatrix}
1 & 0
\end{bmatrix}
\left(
\begin{bmatrix}
s & -1\\
1 & s
\end{bmatrix}^{-1}
\right)
\begin{bmatrix}
0\\
1
\end{bmatrix}
=
\frac{1}{s^2+1}.
\]

I poli sono in
\[
s=\pm j.
\]
Essendo poli sull'asse immaginario, il sistema \textbf{non è BIBO stabile}. Infatti un ingresso sinusoidale in risonanza produce un'uscita non limitata.

Ad esempio, con
\[
u(t)=\sin t,
\]
si ottiene una risposta contenente un termine del tipo
\[
t\cos t,
\]
che cresce in modulo e quindi non resta limitato.

\medskip

\noindent
Questo esempio è molto istruttivo:
\begin{itemize}
    \item il sistema non è asintoticamente stabile internamente;
    \item il sistema non è BIBO stabile esternamente;
    \item le due nozioni però restano concettualmente distinte.
\end{itemize}

\bigskip

\section{Ricapitolazione finale}

Raccogliamo i messaggi centrali della lezione.

\subsection*{1. Realizzazione MIMO}
Per un sistema MIMO si parte dalla matrice di trasferimento \(G(s)\). Per capirne la struttura conviene scriverla come
\[
G(s)=\frac{1}{d(s)}P(s),
\]
studiare la forma di Smith di \(P(s)\) e poi la forma di Smith--McMillan di \(G(s)\).

\subsection*{2. Poli, zeri e ordine minimo}
Dalla forma di Smith--McMillan si leggono:
\begin{itemize}
    \item i poli della matrice di trasferimento;
    \item gli zeri di trasmissione;
    \item il grado di McMillan, cioè la dimensione minima di una realizzazione.
\end{itemize}

\subsection*{3. Zeri di trasmissione}
Uno zero di trasmissione è un valore \(\delta\) tale che esiste un ingresso esponenziale \(e^{\delta t}\) e una opportuna condizione iniziale che rendono l'uscita identicamente nulla.

\subsection*{4. Stabilità interna}
Dipende dagli autovalori di \(A\), quindi dai modi dello stato.

\subsection*{5. Stabilità esterna}
Dipende dai poli della funzione di trasferimento, cioè solo dai modi raggiungibili e osservabili.

\subsection*{6. Differenza tra le due}
In una realizzazione non minima si può avere:
\begin{itemize}
    \item sistema internamente instabile;
    \item ma esternamente stabile,
\end{itemize}
se il modo instabile è nascosto perché non raggiungibile o non osservabile.

\medskip

In altre parole:
\begin{quote}
\emph{la funzione di trasferimento non racconta tutto dello stato interno: racconta solo ciò che è visibile dall'ingresso e dall'uscita.}
\end{quote}